% Part: sets-functions-relations
% Chapter: functions
% Section: isomorphisms

\documentclass[../../../include/open-logic-section]{subfiles}

\begin{document}

\olfileid{sfr}{fun}{iso}
\olsection{Isomorfisme}

\begin{explain}
\emph{Isomorfisme} minangka bijeksi kang njaga struktur himpunan kang
disambungake. Struktur kasebut ditemtokake dening sesambungan kang
lumaku antarane !!{element}s ing himpunan. Gatekna loro himpunan
$X=\{1,2,3\}$ lan $Y=\{4,5,6\}$. Kalorone nduweni struktur
saka relasi penerus, luwih cilik tinimbang, lan luwih gedhe tinimbang.
Isomorfisme antarane loro himpunan kasebut minangka !!a{bijection}
kang njaga struktur kasebut. Dadi fungsi !!a{bijective}
$f \colon X \to Y$ minangka isomorfisme yen $i<j$ yen lan mung yen
$f(i)<f(j)$, $i>j$ yen lan mung yen $f(i)>f(j)$, lan $j$
minangka penerus saka $i$ yen lan mung yen $f(j)$ minangka penerus
saka $f(i)$.
\end{explain}

\begin{defn}[Isomorfisme]
Upamane $U$ minangka pasangan $\langle X, R\rangle$ lan $V$
minangka pasangan $\langle Y, S\rangle$, kanthi $X$ lan $Y$
minangka himpunan, sarta $R$ lan $S$ minangka relasi ing $X$ lan
$Y$ kanthi urut. !!{bijection} $f$ saka $X$ menyang $Y$ minangka
\emph{isomorfisme} saka $U$ menyang $V$ yen lan mung yen njaga
struktur relasional, yaiku kanggo saben $x_{1}$ lan $x_{2}$ ing $X$,
$\tuple{x_1,x_2} \in R$ yen lan mung yen
$\tuple{f(x_1),f(x_2)} \in S$.
\end{defn}

\begin{ex}
Gatekna loro himpunan $X=\{1,2,3\}$ lan $Y=\{4,5,6\}$,
sarta relasi luwih cilik tinimbang lan luwih gedhe tinimbang.
Fungsi $f\colon X \to Y$ kanthi $f(x) = 7-x$ minangka
isomorfisme antarane $\tuple{X,<}$ lan $\tuple{Y,>}$.
\end{ex}

\end{document}
